\documentclass[11pt,a4paper]{article}
\usepackage[margin=23mm,top=25mm,bottom=24mm]{geometry}
\usepackage{fontspec,polyglossia,amsmath,amssymb,booktabs,array,enumitem,xcolor,fancyhdr,xurl,hyperref,microtype}
\setdefaultlanguage{spanish}\setmainfont{Latin Modern Roman}\setsansfont{DejaVu Sans}
\definecolor{ink}{HTML}{153A45}\definecolor{mint}{HTML}{087D78}
\hypersetup{colorlinks=true,urlcolor=mint,linkcolor=ink,pdfauthor={Juan Manuel Goenaga Castro}}
\setlength{\parindent}{0pt}\setlength{\parskip}{6pt}\emergencystretch=3em
\setlength{\headheight}{16pt}\pagestyle{fancy}\fancyhf{}
\fancyhead[L]{\sffamily\footnotesize GoToGym Quantum Research / Music Lab}
\fancyhead[R]{\sffamily\footnotesize M05 / v0.1}
\fancyfoot[L]{\scriptsize Prepublicación metodológica / 22 de octubre de 2026}
\fancyfoot[R]{\small\thepage}\setlist{nosep,leftmargin=1.5em}
\begin{document}
\begin{center}{\sffamily\small\textcolor{mint}{SERIE QAF / TECNOLOGÍA MUSICAL / M05}}\par\vspace{8pt}
{\sffamily\LARGE\bfseries\textcolor{ink}{Validación musical y propiedad intelectual}}\par\vspace{10pt}
\textbf{Juan Manuel Goenaga Castro}\\GoToGym Quantum Research\\22 de octubre de 2026
\end{center}
\colorbox{ink!6}{\parbox{.96\linewidth}{\small\textbf{Estado científico.} Propuesta matemática con prototipo ejecutable. Sin revisión por pares ni validación de preferencia musical. Las extensiones neuronales se identifican como trabajo futuro.}}
\begin{abstract}
Se propone cerrar el ciclo de mejora musical con evaluación pareada, trazabilidad de activos y límites explícitos a las afirmaciones. El prototipo implementa bootstrap por oyente y verificación declarativa de procedencia ligada al hash del archivo. No sustituye una revisión contractual ni aprueba automáticamente una campaña. Se formula un mapa de uso QAF condicionado a evidencia, con opciones de empate o no uso. El resultado permite distinguir software correcto, música preferida y mejora comercial.
\end{abstract}
\section{Por qué un sistema necesita esta capa}
Una nueva fórmula puede generar audio válido sin que el público lo encuentre atractivo. A su vez, una pista preferida puede carecer de permisos para una campaña. La decisión de publicación necesita ambos caminos: calidad y derechos. Los documentos CSRA y MPS/TDVP/RDM suministrados muestran mecanismos experimentales específicos y reconocen límites; no contienen un estudio de preferencia para esta música.

La pregunta es cómo decidir una mejora sin usar el propio algoritmo como juez exclusivo. Se propone comparar contra la música anterior y contra una referencia profesional con licencia. Un productor, editor, ingeniero de mezcla, responsable de marca, investigador de audiencia e interlocutor de derechos aportan criterios diferentes. No se afirma que un programa haya adquirido sus credenciales profesionales.
\section{Variables y diseño}
Cada fila de escucha registra oyente, episodio, valoración del control y del candidato en la misma escala de cero a cien. Actualidad, ajuste de marca y comprensión deben almacenarse como métricas separadas y con una primaria preregistrada. La función incluida procesa una métrica a la vez. Los datos de ejemplo y tests son sintéticos y se rotulan; no se registran como participantes reales.

Orden y etiqueta del método se ocultan. Se conserva voz, imagen y sonoridad; se aleatoriza exposición, se comprueba dispositivo y familiaridad. La repetición de un mismo video no se cuenta como evidencia de muchos episodios independientes.

\newpage
\section{Estimación implementada}
Para oyente $l$, episodio $e$ y diferencia $d_{le}=s^{\mathrm{nuevo}}_{le}-s^{\mathrm{control}}_{le}$:
\begin{equation}
 \bar d_l=\frac{1}{n_l}\sum_ed_{le},\qquad
 \hat\Delta=\frac{1}{L}\sum_l\bar d_l.
\end{equation}
Cada oyente tiene igual peso. El bootstrap remuestrea oyentes completos y obtiene percentiles 2.5 y 97.5 de la media de sus diferencias. Es un intervalo condicionado a los episodios evaluados; no estima incertidumbre entre campañas ni convierte escalas ordinales en evidencia causal universal.

Se rechazan filas duplicadas del mismo par oyente-episodio. Si hay repeticiones deben agregarse antes, con regla registrada. El mínimo de 20 oyentes y cinco episodios es una compuerta de piloto del software, no un cálculo de potencia. El estado favorable se llama \texttt{pilot\_signal}; jamás ``eficacia probada''. El tamaño confirmatorio se calculará con variabilidad de un piloto real.
\section{Extensión confirmatoria propuesta}
Un modelo cruzado permitiría separar oyente y episodio:
\begin{equation}
 d_{le}=\mu+u_l+v_e+\varepsilon_{le},\quad
 u_l\sim\mathcal N(0,\sigma_l^2),\quad v_e\sim\mathcal N(0,\sigma_e^2).
\end{equation}
Su adecuación dependerá de distribución y escala. No está ajustado en v0.1. La prueba de ventas debe ser otro experimento con asignación, tráfico y denominadores verificables; una preferencia auditiva no prueba más conversiones.
\section{Regiones de decisión QAF propuestas}
Para ganancia empírica $\Delta$, intervalo $[a,b]$, costo $C$ y presupuesto $C_{\max}$: QAF-positive solo si el criterio preregistrado se supera fuera de calibración y dentro de presupuesto; classical-positive si el control lo supera; non-identifiable si no se separan; computational-breakdown si el costo impide operar. Son regiones para un benchmark futuro, no etiquetas adjudicadas al prototipo. Los algoritmos actuales de selección y alineación coinciden con sus controles clásicos de igual información.

\newpage
\section{Contrato de procedencia implementado}
\texttt{verify\_asset} exige hash SHA-256, origen, identificador de licencia, referencia a evidencia de derechos, revisor, fecha de revisión y vigencia. Requiere declaraciones explícitas de uso comercial y sincronización. Calcula el hash de los bytes y comprueba fecha de corte dentro de la vigencia declarada:
\begin{equation}
 H(\mathrm{archivo})=h_{\mathrm{registrado}},\qquad
 t_{\mathrm{revisión}}\leq t_{\mathrm{uso}}\leq t_{\mathrm{vigencia}}.
\end{equation}
Estas igualdades vinculan evidencia declarada y activo; no autentican un contrato ni prueban titularidad. Los documentos de permisos deben estar en almacenamiento privado. No se carga una voz del usuario al repositorio para evaluar música.

La licencia de código, de pesos neuronales, de grabación, de composición y de sincronización son objetos distintos. El código DSP original evita incorporar muestras descargadas, pero no acredita automáticamente exclusividad o registrabilidad de toda secuencia generada. Un modelo neuronal futuro requerirá dataset card, versión de pesos y trazabilidad de entrenamiento.
\section{Pruebas ejecutadas y fallos esperados}
El test estadístico crea cuatro oyentes y dos episodios con ganancia constante de 20 puntos; confirma esa media pero conserva evidencia insuficiente por tamaño. Un duplicado se rechaza. El test de procedencia acepta un fixture completo, rechaza permiso de sincronización falso y rechaza bytes modificados después de registrar el hash.

La suite específica contiene doce controles de ingeniería, incluyendo selección, sincronía exacta, síntesis reproducible y mezcla FFmpeg. Su éxito no equivale a doce oyentes ni a doce videos preferidos. Los informes de validación conservan comandos, condición sintética y ausencia de estudio humano.
\section{Protocolo de publicación}
Primero audicionar tres bocetos; después comparar el seleccionado con productor o pista profesional autorizada; luego mezclar con voz real y revisar el archivo codificado. Conservar perdedores y parámetros. Cualquier cambio posterior de activo invalida su aprobación hasta nueva revisión. El checker nunca escribe una aprobación artística verdadera ni autoriza una campaña por la existencia de un hash.

\newpage
\paragraph{Propiedad intelectual y fronteras.}
La contribución propuesta por Juan Manuel Goenaga Castro consiste en integrar contexto de marca, control compositivo y evidencia de decisión en el sistema audiovisual GoToGym. Se conservan las atribuciones del QAF y del antecedente técnico de Antoni Bancells. DSP, teoría musical, programación dinámica, bootstrap, NumPy y FFmpeg no se presentan como invenciones propias. No se determina patentabilidad ni exclusividad. Los originales privados, voces, contratos y referencias comerciales no se redistribuyen. Una licencia del código no equivale a permiso sobre pesos, muestras, composición, máster o sincronización.
\section{Aprendizajes y criterios de parada}
La evidencia debe poder decir que el nuevo sistema no mejora. Si un track profesional produce mayor preferencia a menor costo, usarlo para la campaña y conservar el laboratorio como desarrollo. Si falta un derecho o un dato de evaluación, declarar ese hueco sin sustituirlo por un índice agregado.

El capital intelectual documentable incluye contratos entre módulos, modelos matemáticos explícitos, decisiones de marca, diseños de evaluación y código reproducible. No incluye apropiarse de estándares de sonoridad, teoría musical ni resultados de terceros. La publicación propuesta conserva esas fronteras y no presenta la propiedad intelectual como garantía de calidad.
\section{Conclusión}
Se entrega un mecanismo para aprender de la audiencia y controlar qué evidencia acompaña cada archivo. La mejora musical se podrá afirmar cuando una comparación independiente la sostenga. Mientras tanto, la plataforma entrega bocetos, controles y un camino verificable hacia esa prueba.
\section*{Referencias y trazabilidad}
\begin{enumerate}\small
\item Goenaga Castro, J. M. QAF (2025); \emph{Ventaja estructural CSRA} y \emph{Motor completo MPS/TDVP/RDM} (2026). Manuscritos aportados, consultados; inventario SHA-256 en \texttt{validation.json}. Su evidencia no se transfiere a música.
\item Bancells, A. (2026). Antecedente técnico MPS/TDVP/QUBO, atribución conservada en la serie audiovisual del repositorio. No se reejecuta su motor en esta extensión.
\item Copet et al. (2023). \emph{Simple and Controllable Music Generation}. NeurIPS 36. \url{https://arxiv.org/abs/2306.05284}.
\item ITU-R BS.1770-5 y BS.1534-3. Sonoridad y evaluación subjetiva. \url{https://www.itu.int/rec/R-REC-BS.1770} y \url{https://www.itu.int/rec/R-REC-BS.1534}.
\item FFmpeg. Filtros \texttt{loudnorm} y \texttt{sidechaincompress}. \url{https://ffmpeg.org/ffmpeg-filters.html}.
\item Meta. Licencia de pesos AudioCraft. \url{https://github.com/facebookresearch/audiocraft/blob/main/LICENSE_weights}. Consulta de fuentes externas: 22/10/2026.
\end{enumerate}

\end{document}
